% !TeX program = pdflatex
\documentclass[11pt,a4paper]{article}
\usepackage[T1]{fontenc}
\usepackage[utf8]{inputenc}
\usepackage{lmodern,amsmath,amssymb,amsthm,mathtools,microtype}
\usepackage[margin=25mm]{geometry}
\usepackage{booktabs,enumitem,xurl}
\usepackage[hidelinks]{hyperref}
\setlength{\emergencystretch}{2em}
\setlist{nosep,leftmargin=*}
\newtheorem{theorem}{Theorem}[section]
\newtheorem{proposition}[theorem]{Proposition}
\newtheorem{lemma}[theorem]{Lemma}
\newtheorem{corollary}[theorem]{Corollary}
\theoremstyle{definition}
\newtheorem{definition}[theorem]{Definition}
\theoremstyle{remark}
\newtheorem{remark}[theorem]{Remark}
\newcommand{\Ric}{\operatorname{Ric}}
\newcommand{\CP}{\mathbb{CP}}
\newcommand{\Met}{\operatorname{Met}}
\newcommand{\PD}{\operatorname{PD}}
\newcommand{\Span}{\operatorname{span}}
\newcommand{\pair}[2]{\langle #1,#2\rangle}
\title{Global period constraints for the Q3 minimizing problem\\[3pt]
\large Marked negative lattices, cyclic bubbles, and the Godeaux blowup}
\author{Qiushi research continuation}
\date{8 October 2026}
\begin{document}
\maketitle
\begin{abstract}
We keep the smooth-attainment problem for the trace-free Ricci functional on
one blowup of the classical Godeaux surface. A marked negative lattice whose
self-dual periods vanish forces the pullback canonical class to have a larger
self-dual square. Combining this fact with the classical monopole curvature
inequality gives a global energy restriction on a minimizing sequence, without
an energy identity or a monopole extension across the limiting orbifold.
In particular, two distinct degree-one $(-3)$ classes of the classical
Godeaux configuration cannot both enter such a degeneration below $8\pi^2$.
A single such class requires total energy at least $64\pi^2/9$, including
when the mass-replacement eigenvalue is zero. A further corollary restricts
canonically marked negative cyclic chains. The results do not establish smooth
attainment: canonical-orthogonal bubbles, a single permitted charged chain,
and other degeneration channels remain.
\end{abstract}

\section{Fixed problem and the classical input}
Let $X$ be the classical Godeaux surface and choose the blowup point outside
its fifteen degree-one rational curves. Put
\[
 M=X\#\overline{\CP}^{\,2},\qquad
 K=\pi^*c_1(K_X),\qquad E^2=-1,\quad K\cdot E=0,\quad K^2=1.
\]
All cohomology calculations are in $H^2(M;\mathbb R)$, so torsion is omitted.
For a metric $g$ let
\[
 P(g)=\int_M|\Ric_g^\circ|^2dV_g,\qquad
 p_*(M)=\inf_{g\in\Met(M)}P(g),\qquad T=2\chi(M)+3\tau(M)=0.
\]
The supplied Q3 construction gives the strict upper bound
\begin{equation}\label{eq:upper}
 \frac{16\pi^2}{3}\le p_*(M)<8\pi^2.
\end{equation}
The strict inequality is an input from the general negative-K\"ahler--Einstein
insertion construction in the supplied October 8 manuscripts. It is not
reproved here. All conformal classes on this $M$ are Yamabe-negative; a
smooth $P$-minimizer would have constant negative scalar curvature and,
if Bach-flat, would be Einstein. The positive lower bound excludes this
last alternative. Thus smooth attainment on this fixed $M$ would answer Q3.
These are the same reductions used in the supplied mass-variation note
\cite{Source,MassNote}.

We use tensor norms for Ricci curvature and the Hilbert--Schmidt
curvature-operator norm for Weyl curvature. The classical estimate for an
actual monopole class $a$ is
\[
 \frac1{4\pi^2}\int_M\left(\frac{s^2}{24}+2|W^+|^2\right)dV_g
 \ge\frac23\|a_g^+\|^2.
\]
Together with Chern--Weil this gives
\begin{equation}\label{eq:classical}
 P(g)\ge\frac{16\pi^2}{3}\|a_g^+\|^2-8\pi^2T.
\end{equation}
The usual blowup supply selects an actual sign class with self-dual
square at least $\|K_g^+\|^2$; this is the projection argument in
LeBrun's Proposition 3.2, not just its final numerical bound
\cite{LeBrun}. In our case the classes are $K\pm E$, up to an overall sign.
Consequently the only geometric input needed for the main argument is
\begin{equation}\label{eq:Kbound}
 \boxed{P(g)\ge\frac{16\pi^2}{3}\|K_g^+\|^2.}
\end{equation}
This keeps the actual-supply and, where relevant, chamber conventions of that
classical argument. It does not require FD, the O10 endpoint theorem, or any
new equation on the limiting orbifold.

\section{A marked lattice changes the global projection budget}
Let $C_1,\ldots,C_r\in H^2(M;\mathbb R)$ be linearly independent marked
classes with negative-definite matrix
\[
 Q=(C_i\cdot C_j),\qquad d_i=K\cdot C_i.
\]
Define
\begin{equation}\label{eq:beta}
 v=Q^{-1}d,\qquad L=\sum_i v_iC_i,\qquad
 \beta_Q(K)=-d^{\mathsf T}Q^{-1}d\ge0,\qquad K_\perp=K-L.
\end{equation}
These are intersection-form projections, not metric-dependent Hodge
projections. Directly,
\begin{equation}\label{eq:perp}
 K_\perp\cdot C_i=0,\qquad L^2=-\beta_Q(K),\qquad
 K_\perp^2=1+\beta_Q(K).
\end{equation}

\begin{theorem}[Finite-metric period bound]\label{thm:finite}
For every smooth metric on the fixed $M$,
\begin{equation}\label{eq:finite}
 \boxed{
 P(g)\ge\frac{16\pi^2}{3}
 \left(\sqrt{1+\beta_Q(K)}-\|L_g^+\|\right)_+^2.
 }
\end{equation}
In particular, if $g_j$ is any sequence for which
$\|(C_i)^+_{g_j}\|\to0$ for every $i$, then
\begin{equation}\label{eq:limit}
 \boxed{\liminf_jP(g_j)\ge
 \frac{16\pi^2}{3}\left(1+\beta_Q(K)\right).}
\end{equation}
\end{theorem}
\begin{proof}
Hodge decomposition gives
\[
 \|(K_\perp)^+_g\|^2
 =K_\perp^2+\|(K_\perp)^-_g\|^2\ge1+\beta_Q(K).
\]
Since $K^+=(K_\perp)^++L^+$, the triangle inequality yields
\[
 \|K_g^+\|\ge
 \left(\sqrt{1+\beta_Q(K)}-\|L_g^+\|\right)_+.
\]
Use \eqref{eq:Kbound}. For the second conclusion,
$\|L_{g_j}^+\|\le\sum_i|v_i|\|(C_i)^+_{g_j}\|\to0$.
\end{proof}

\begin{corollary}[Restriction on a minimizing degeneration]\label{cor:budget}
If the periods in Theorem~\ref{thm:finite} vanish along a minimizing
sequence, then
\begin{equation}\label{eq:budget}
 \boxed{\beta_Q(K)\le\frac{3p_*(M)}{16\pi^2}-1<\frac12.}
\end{equation}
Equivalently, any marked negative lattice with $\beta_Q(K)\ge1/2$
is excluded from this kind of degeneration.
\end{corollary}
\begin{proof}
Combine \eqref{eq:limit} with \eqref{eq:upper}.
\end{proof}

The estimate is on the original closed manifold. In particular, it is
not obtained by adding a global lower bound to the energy of a bubble:
that would in general count some curvature twice. No energy equality,
compactness of the exterior, or existence of a bulk Einstein metric
occurs in its proof.

\subsection{The variant retaining the complete sign pair}
When both sign classes are available for the metric under consideration,
the same argument keeps possible intersections with the new exceptional
class. Set $d_E=(E\cdot C_i)_i$ and, for $|t|\le1$, let
\[
 B_Q(t)=1-t^2-(d+td_E)^{\mathsf T}Q^{-1}(d+td_E).
\]
Since $K+tE$ is a convex combination of $K+E$ and $K-E$,
\[
 \|(K+tE)^+\|^2\le
 \max\{\|(K+E)^+\|^2,\|(K-E)^+\|^2\}.
\]
Projecting $K+tE$ off the negative lattice as above gives, when its
marked self-dual periods vanish,
\begin{equation}\label{eq:fullsign}
 \liminf_jP(g_j)\ge
 \frac{16\pi^2}{3}\max_{-1\le t\le1}B_Q(t).
\end{equation}
Here the maximum is of an explicit quadratic polynomial on a compact
interval. If $d_E=0$, its maximum is $1+\beta_Q(K)$, so it reduces to
\eqref{eq:limit}. The displayed sign-pair hypothesis should be retained
when \eqref{eq:fullsign} is applied outside the present example.

\section{From a genuine ALE region to vanishing self-dual periods}
We now give a sufficient geometric hypothesis for the periods above.
It concerns an identified marked ALE limit; it does not assert that every
minimizing sequence has such a limit.

\begin{lemma}[Compactly supported almost anti-self-dual representatives]
Let $(N,h)$ be a smooth oriented complete ALE four-manifold with
negative-definite compact intersection form. Every compactly supported
cohomology class has representatives $\omega_R$, closed and supported
in a compact truncation of the end, with
\[
 \|\omega_R^+\|_{L^2(h)}\longrightarrow0.
\]
\end{lemma}
\begin{proof}
On an ALE four-manifold, $H^2_c(N;\mathbb R)\to H^2(N;\mathbb R)$ is an
isomorphism, and its classes have unique $L^2$ harmonic representatives
\cite[Lemma 1.7]{HL}. The middle-dimensional $L^2$ pairing agrees with
the compact-support intersection form. A nonzero self-dual harmonic form
would give a positive square, so negative definiteness implies that all
these harmonic representatives $\Omega$ are anti-self-dual.

On an annulus $A_R=\{R<r<2R\}$, $H^2(A_R;\mathbb R)=0$ because the link
is a spherical space form. Choose $\alpha_R$ with $d\alpha_R=\Omega$.
The fixed-annulus Hodge estimate, after scaling and using the ALE
convergence of the metric, permits the choice
\[
 \|\alpha_R\|_{L^2(A_R)}\le C R\|\Omega\|_{L^2(A_R)}.
\]
One may enlarge the annulus slightly for the support of a cutoff.
Let $\chi_R$ be zero near its inner boundary and one near its outer
boundary, with $|d\chi_R|\le C/R$. Put $\omega_R=\Omega$ inside,
$\omega_R=\Omega-d(\chi_R\alpha_R)$ on the annulus, and zero outside.
This is a smooth compactly supported closed representative of the same
class. Its self-dual part is confined to the annulus and satisfies
\[
 \|\omega_R^+\|_2
 \le C\|\Omega\|_{L^2(A_R)}\longrightarrow0.
\]
The last limit follows from the finite $L^2$ norm. This also shows why
no uniform curvature bound on the exterior of a bubble is needed.
\end{proof}

\begin{proposition}[Marked ALE concentration implies the required periods]
Suppose that a fixed finite collection of marked compact classes in $M$
is carried by smooth negative-definite ALE limits, with smooth rescaled
convergence on every fixed compact truncation and with the compact-support
classes identified with the specified $C_i$. Then
\begin{equation}\label{eq:vanish}
 \|(C_i)^+_{g_j}\|\longrightarrow0.
\end{equation}
The assertion applies in particular to a genuine marked smooth
$\mathcal O(-3)$ ALE bubble. It requires neither an energy equality nor
regularity of the entire limiting bulk.
\end{proposition}
\begin{proof}
Fix $R$ and pull $\omega_R$ from the lemma into $M$ through the marked
convergence map; extend by zero. The result is a closed representative
of $C_i$. On a closed manifold the norm of the harmonic self-dual
projection is at most the self-dual norm of any closed representative.
Moreover, the $L^2$ norm of a two-form and the Hodge star on two-forms
are unchanged by constant metric scaling in dimension four. Consequently,
\[
 \limsup_j\|(C_i)^+_{g_j}\|\le\|\omega_R^+\|_{L^2(h)}.
\]
Now let $R\to\infty$.
\end{proof}

For a tree with orbifold intermediate levels, the relevant markings and
period limits must still be obtained on the original manifold. The
proposition covers smooth ALE pieces or any other region for which
\eqref{eq:vanish} has been established; it does not turn an arbitrary
orbifold or collapsed limit into such a piece.

\section{Consequences for the classical Godeaux configuration}
\subsection{One or two \texorpdfstring{$(-3)$}{(-3)} classes}
A degree-one rational curve $C$ in $X$ has
\[
 C^2=-3,\qquad K\cdot C=1,\qquad E\cdot C=0.
\]
For its rank-one lattice, $Q=[-3]$ and $\beta_Q(K)=1/3$. Hence:
\begin{corollary}\label{cor:one}
If this marked class enters a degeneration covered by
\eqref{eq:vanish}, then
\begin{equation}\label{eq:one}
 \liminf_jP(g_j)\ge\frac{64\pi^2}{9}
 =\frac89(8\pi^2).
\end{equation}
This includes the first-order zero-mode case of the mass-replacement
argument.
\end{corollary}
For two disjoint such classes the charge is $2/3$, so instead
\begin{equation}\label{eq:two}
 \liminf_jP(g_j)\ge\frac{80\pi^2}{9}
 =\frac{10}{9}(8\pi^2)>8\pi^2.
\end{equation}
This contradicts the inherited strict upper bound for a minimizing
sequence. Although two separate scalar-flat K\"ahler $\mathcal O(-3)$
fillings have combined local $P$-energy only $16\pi^2/3$, their
simultaneous marked degeneration is globally excluded by
\eqref{eq:two}.

\subsection{At most one of the fifteen original degree-one lines}
The classical Godeaux surface has fifteen degree-one rational curves
$L_i^m$, with $m\in\{-1,0,1\}$ and $i\in\mathbb Z/5$. Their intersections
are computed in \cite[Proposition 4.3]{Godeaux}. Distinct curves have
intersection $s\in\{0,1,2,5\}$; each has square $-3$ and canonical
degree one. If $s<3$, the pair matrix is negative definite and
\[
 \beta_Q(K)=\frac2{3-s}.
\]
Thus the lower bounds, normalized by $8\pi^2$, are
\[
\begin{array}{c|ccc}
s&0&1&2\\\hline
\beta_Q(K)&2/3&1&2\\
\displaystyle\liminf P/(8\pi^2)&10/9&4/3&2
\end{array}
\]
and are all incompatible with a minimizing sequence. If $s=5$, then
$(C+D)^2=4$, so $\|(C+D)^+\|\ge2$ for every metric; the two self-dual
periods cannot both tend to zero even without an energy estimate.

\begin{corollary}\label{cor:15}
Along a minimizing sequence on this marked Godeaux blowup, at most one
of the fifteen original degree-one rational-curve classes can have
vanishing self-dual period. In particular at most one can enter the
negative-definite ALE degenerations described in the previous section.
\end{corollary}
The statement refers to these fixed classes, not to a classification of
all embedded spheres or all possible markings on $M$.

\section{Canonically marked cyclic chains}
For a proposed cyclic plumbing let
\[
 Q_{ii}=-n_i,\quad Q_{i,i+1}=Q_{i+1,i}=1,\quad n_i\ge2,
 \qquad K\cdot C_i=n_i-2,
 \qquad E\cdot C_i=0.
\]
The pairings are part of the marking hypothesis. They hold by adjunction
for a rational curve configuration inherited from the unblown surface;
they must not be assigned to an arbitrary smooth bubble without proof.
Assume these periods tend to zero along a minimizing sequence.

\begin{theorem}[Restriction on cyclic chain data]\label{thm:chain}
Every connected chain with a nonzero canonical pairing must have the form
\[
 [-3,\underbrace{-2,\ldots,-2}_{r-1}]
\]
or its reversal. For this chain,
\begin{equation}\label{eq:chain}
 \beta_Q(K)=\frac{r}{2r+1},\qquad
 \liminf_j\frac{P(g_j)}{8\pi^2}\ge1-\frac1{6r+3}.
\end{equation}
At most one mutually orthogonal component can have a nonzero canonical
pairing. Chains of $(-2)$ classes are not excluded by this theorem.
\end{theorem}
\begin{proof}
Write $A=-Q$, so $A$ has diagonal $n_i$ and off-diagonal $-1$, and set
$d_i=n_i-2\ge0$. Its charge has the variational expression
\[
 d^{\mathsf T}A^{-1}d
 =\max_z\{2d^{\mathsf T}z-z^{\mathsf T}Az\}.
\]
Testing with $z_i=d_i/n_i$ gives
\[
 \beta_Q(K)\ge\sum_i\frac{(n_i-2)^2}{n_i}.
\]
If some $n_i\ge4$, the right side is at least one. If two weights are
at least three, it is at least $2/3$. Both contradict
\eqref{eq:budget}. Hence there is exactly one weight three and all other
weights are two.

Let the two $(-2)$ tails have lengths $l$ and $r'$. The endpoint diagonal
entry of the inverse $A_l$ Cartan matrix is $l/(l+1)$, as follows directly
from the tridiagonal recurrence. Its Schur complement gives
\[
 \beta_Q(K)=
 \frac1{3-l/(l+1)-r'/(r'+1)}.
\]
If both tails are nonempty, this is at least $1/2$, again excluded. The
three must therefore be an end vertex. With total length $r$ the charge
is $r/(2r+1)$, and Theorem~\ref{thm:finite} gives \eqref{eq:chain}.
For orthogonal components the charges add. Each nonzero component has
charge at least $1/3$, so two would violate \eqref{eq:budget}.
\end{proof}

For lengths $r=1,2,3,4$, the normalized thresholds are respectively
$8/9,14/15,20/21,26/27$. A verified upper bound $p_*/(8\pi^2)\le U<1$
restricts any such surviving chain to
\[
 6r+3\le\frac1{1-U}.
\]
This is a restriction on specified marked cyclic data, not an assertion
that all minimizing bubbles are cyclic or canonically marked.

\section{Relation to the mass zero mode and to attainment}
The supplied mass-variation result uses a different input: a smooth
$P$-critical orbifold point, marked matching, and either a global energy
identity or the stated local lower comparison. It gives
\[
 P(g_{\lambda,t})=P(g_{\rm orb})+\frac{8\pi^2}{3}
 -24\pi^2m(g_t)\mu\lambda^2+o(\lambda^2),
\]
where positive and negative masses are available on the same filling.
Under those hypotheses, a minimizing degeneration requires $\mu=0$
\cite{MassNote}. The present period argument neither changes this
expansion nor uses its coefficient. It imposes an independent condition
on the same marked degeneration:
\[
 \boxed{\mu=0,\qquad p_*(M)\ge64\pi^2/9}
\]
when both sets of hypotheses apply. Two of the original $(-3)$ classes
cannot occur together, regardless of their local values of $\mu$.

The single surviving charged case is not excluded because
$64\pi^2/9<8\pi^2$. If a future strict upper bound below $64\pi^2/9$ is
proved on this same $M$, all rank-one canonical-degree-one $(-3)$
degenerations satisfying the period hypothesis will then be excluded.
No such upper bound is proved here.

If $K$ pairs trivially with a negative lattice, its charge in
\eqref{eq:beta} is zero. Thus the method does not exclude canonical-orthogonal
ADE lattices or bubbles with no real second cohomology. It also does not
establish marked convergence for every minimizing sequence or rule out
partial collapse. The separate zero-mode calculation found in the
companion research records likewise retains an undetermined full
fourth-order gluing coefficient \cite{ZeroMode}; a negative inner
coefficient alone cannot remove this final case.

The current fixed task is unchanged: obtain smooth global attainment on
$M$. The increment proved here is a constraint on which actual marked
negative-lattice degenerations a minimizing sequence could have. It
uses the original closed-manifold monopole inequalities throughout,
rather than presupposing an extension of those inequalities to an
unknown orbifold bulk.

\begin{thebibliography}{9}
\bibitem{LeBrun} C. LeBrun, \emph{Ricci Curvature, Minimal Volumes, and
Seiberg--Witten Theory}, Invent. Math. \textbf{145} (2001), 279--316.
Propositions 3.1--3.2 and Corollary 3.4.
\url{https://arxiv.org/abs/math/0003068}.
\bibitem{HL} H.-J. Hein and C. LeBrun, \emph{Mass in K\"ahler Geometry},
Comm. Math. Phys. \textbf{347} (2016), 183--221. Lemma 1.7.
\url{https://arxiv.org/abs/1507.08885}.
\bibitem{Godeaux} C. B\"ohning, H.-C. Graf von Bothmer and P. Sosna,
\emph{On the derived category of the classical Godeaux surface},
Adv. Math. \textbf{243} (2013), 203--231. Section 4, Proposition 4.3.
\url{https://arxiv.org/abs/1206.1830}.
\bibitem{Source} Supplied project manuscripts,
\emph{Four Public Problems: Overnight Results, Proof Mechanisms,
Remaining Connections, and Questions for Further Research}, and
\emph{Mathematical Source Companion}, 8 October 2026.
The strict upper construction and the Q3 variational reduction are inputs.
\bibitem{MassNote} Supplied continuation,
\emph{Q3: Mass Variation and a Cyclic Degeneration Obstruction},
8 October 2026, Sections 2--5.
\bibitem{ZeroMode} Related continuation,
\emph{Q3: The Zero Mode and the Proposed Horikawa Candidate},
8 October 2026, Sections 4--6.
\end{thebibliography}
\end{document}
